\documentclass[11pt]{article}
\usepackage[T1]{fontenc}
\usepackage{lmodern,microtype,amsmath,amssymb,amsthm,mathtools,tikz-cd}
\usepackage[margin=1in]{geometry}
\usepackage{xurl,hyperref}
\hypersetup{colorlinks=true,linkcolor=blue,citecolor=blue,urlcolor=blue,
pdftitle={Syntax and deductively closed theories -- Version 16},pdfauthor={Florian Lengyel}}
\newtheorem{theorem}{Theorem}[section]
\newtheorem{proposition}[theorem]{Proposition}
\newtheorem{lemma}[theorem]{Lemma}
\theoremstyle{definition}\newtheorem{definition}[theorem]{Definition}
\newcommand{\N}{\mathbb N}
\newcommand{\Crec}{\mathcal C_{\mathrm{rec}}}
\newcommand{\Da}{\mathcal D_a}
\newcommand{\Tm}{\mathsf{Tm}}
\newcommand{\Fm}{\mathsf{Fm}}
\newcommand{\FV}{\operatorname{FV}}
\newcommand{\Two}{\mathbf 2}
\newcommand{\Adm}{\mathsf{Adm}}
\newcommand{\im}{\operatorname{im}}
\newcommand{\ran}{\operatorname{ran}}
\newcommand{\Tot}{\operatorname{Tot}}
\newcommand{\Split}{\operatorname{Split}}
\newcommand{\Par}{\operatorname{Par}}
\newcommand{\Mono}{\operatorname{Mono}}
\newcommand{\List}{\mathsf L}
\newcommand{\nil}{\mathsf{nil}}
\newcommand{\cons}{\mathsf{cons}}
\newcommand{\mem}{\mathsf{mem}}
\newcommand{\del}{\mathsf{del}}
\newcommand{\ok}{\mathsf{ok}}
\newcommand{\fail}{\mathsf{fail}}
\begin{document}
\appendix
\section{Syntax and deductively closed theories}
\label{app:syntax-theories}
\noindent\textit{Version 16, 31 October 2026.}\par\medskip

We define first-order syntax by categorical structure and construct
the theory generated by an axiom subobject. The construction uses
parameterized initial algebras for syntax and lists, an initial algebra
of derivations, and regular images.

\subsection{Categorical hypotheses and notation}
\label{app:categorical-hypotheses}

Let $\mathcal D$ be a regular category with finite coproducts.
Thus $\mathcal D$ has finite limits, every morphism factors as a
regular epimorphism followed by a monomorphism, and regular
epimorphisms are stable under pullback. A regular epimorphism is a
coequalizer; the monic part of the factorization is the image.
For the construction of first-order syntax, assume also that finite
coproducts are disjoint and stable under pullback. In particular,
a morphism into $X+Y$ decomposes its source as the coproduct of its
two inverse images. These additional coproduct hypotheses are not
needed for Theorem~\ref{thm:least-theorem-subobject}.

Composition is $gf=g\circ f$. Finite products and coproducts are
chosen and parenthesized to the left. Rebracketing uses the canonical
associativity isomorphisms determined by their projections or injections;
terminal factors are removed by the canonical product unit isomorphisms.
For $X\in\mathcal D$, write $X^0=1$, $X^1=X$, and
$X^{n+1}=X^n\times X$ for $n\geq1$.
For $f:X\to Y$, the morphism $f^n:X^n\to Y^n$ is the corresponding
product morphism, with $f^0=1_1$ and $f^1=f$.
Thus powers of objects denote finite products.

Put $\Two=1+1$, with injections
$\mathbf f,\mathbf t:1\to\Two$. Pulling this coproduct back along
$\pi_1:\Two\times\Two\to\Two$ decomposes $\Two\times\Two$ into two copies of
$\Two$. Decomposing each copy as $1+1$ gives the isomorphism
\[
 \gamma=[\langle\mathbf f,\mathbf f\rangle,
          \langle\mathbf f,\mathbf t\rangle,
          \langle\mathbf t,\mathbf f\rangle,
          \langle\mathbf t,\mathbf t\rangle]:
          1+1+1+1\xrightarrow{\cong}\Two\times\Two.
\]
Define the Boolean morphisms by
\[
 \neg_2=[\mathbf t,\mathbf f]:\Two\to\Two,\qquad
 \wedge_2=[\mathbf f,\mathbf f,\mathbf f,\mathbf t]\gamma^{-1},
 \qquad
 \vee_2=[\mathbf f,\mathbf t,\mathbf t,\mathbf t]\gamma^{-1},
\]
where $\wedge_2,\vee_2:\Two\times\Two\to\Two$.
Together with $\mathbf f$ and $\mathbf t$, these morphisms make
$\Two$ a Boolean algebra object in $\mathcal D$.
To verify the axioms, distributivity expresses each finite power
$\Two^n$ as a coproduct of $2^n$ copies of $1$.
Its coproduct injections are jointly epimorphic, so each required
equation can be checked after precomposition with these injections.
The displayed copairings give the same composite on each summand
for the two sides of every Boolean algebra axiom. For example,
\[
 \wedge_2\langle1_{\Two},\neg_2\rangle=\mathbf f!_{\Two},
 \qquad
 \vee_2\langle1_{\Two},\neg_2\rangle=\mathbf t!_{\Two}:
 \Two\to\Two,
\]
as follows by precomposing with $\mathbf f$ and $\mathbf t$.

Let $p:X\to\Two$. Write $[p=t]\hookrightarrow X$ for its
pullback along $\mathbf t$; write $[p=f]\hookrightarrow X$ for its
pullback along $\mathbf f$.
A conditional with test $p$ and branches $f,g:X\to Y$ is the
morphism obtained by copairing their restrictions to these two
pullbacks.

We use variables to abbreviate projections from a displayed product.
An expression built from them denotes a composite of morphisms on
that product. An equation between such expressions is an equality
of morphisms with that common source. A morphism $c:1\to Y$ used
on a source $X$ denotes $c!_X:X\to Y$.

Fix an object $V\in\operatorname{Ob}(\mathcal D)$ and a morphism
$\epsilon:V\times V\to\Two$ in $\mathcal D$ such that the square
\[
\begin{tikzcd}
 V\arrow[r,"{!_V}"]\arrow[d,hook,"\Delta_V"'] &
 1\arrow[d,hook,"\mathbf t"]\\
 V\times V\arrow[r,"\epsilon"'] & \Two
\end{tikzcd}
\]
is a pullback. Thus $V$ is equipped with decidable equality:
its diagonal is a complemented subobject of $V\times V$, with
complement given by the pullback of $\mathbf f:1\to\Two$
along $\epsilon$.
For $0\leq k\leq d$, fix objects
$\Sigma_k,\Pi_k\in\operatorname{Ob}(\mathcal D)$ to index the
$k$-ary function and relation constructors, respectively.
Here $d$ is a fixed nonnegative integer. Equality is a logical symbol.

For an endofunctor on the $r$-fold product category $\mathcal D^r$
defined using finite products, coproducts, and specified constant objects, its version on
$(\mathcal D/P)^r$ is obtained
by replacing each constant $X$ by $P\times X\to P$ and using
products and coproducts in that slice. An initial algebra is
\emph{stable under parameters} if, for every $P\in\mathcal D$,
its pullback along $P\to1$ is initial for this slice endofunctor.

\subsection{Terms, formulas, and free variables}
\label{app:first-order-syntax}

Define $K:\mathcal D^2\to\mathcal D^2$ by
$K(X,Y)=\bigl(K_1(X,Y),K_2(X,Y)\bigr)$, where
\begin{align}
 K_1(X,Y)&=V+\coprod_{k=0}^{d}\Sigma_k\times X^k,
                                                        \label{eq:term-functor}\\
 K_2(X,Y)&=X^2+\coprod_{k=0}^{d}\Pi_k\times X^k
                  +Y+Y^2+V\times Y.                    \label{eq:formula-functor}
\end{align}
On a morphism $(f,g):(X,Y)\to(X',Y')$, use
$1_V$, $1_{\Sigma_k}\times f^k$, $f^2$,
$1_{\Pi_k}\times f^k$, $g$, $g^2$, and $1_V\times g$ on
the respective summands. Product and coproduct laws give the
functor laws.

\begin{definition}\label{def:categorical-syntax}
A raw named-variable syntax object is an initial $K$-algebra
$\iota:K(\Tm,\Fm)\to(\Tm,\Fm)$ stable under parameters.
We also require a list object $W=\List V$, given by an initial
algebra $[\nil,\cons]:1+V\times W\to W$ stable under parameters.
\end{definition}

The syntax constructors are the following morphisms in $\mathcal D$:
\begin{gather*}
 \eta:V\to\Tm,\qquad f_k:\Sigma_k\times\Tm^k\to\Tm,\\
 {\mathsf{eq}}:\Tm^2\to\Fm,\qquad
 r_k:\Pi_k\times\Tm^k\to\Fm,\\
 \neg:\Fm\to\Fm,\qquad
 {\Rightarrow}:\Fm^2\to\Fm,\qquad
 {\forall}:V\times\Fm\to\Fm.
\end{gather*}
The morphism $\eta$ is the variable constructor.
The quantifier constructor $\forall$ is indexed
by $V$; its binding behaviour is specified by the free-variable and
substitution morphisms constructed below.
For every $K$-algebra $\alpha:K(X,Y)\to(X,Y)$ there is a unique
$(u,v):(\Tm,\Fm)\to(X,Y)$ making the square
\[
\begin{tikzcd}[column sep=large]
 K(\Tm,\Fm)\arrow[r,"\iota"]\arrow[d,"{K(u,v)}"'] &
 (\Tm,\Fm)\arrow[d,"{(u,v)}"]\\
 K(X,Y)\arrow[r,"\alpha"'] & (X,Y)
\end{tikzcd}
\]
commute. The same property holds in every slice with parameters.
Induction means applying this property to a monic subalgebra:
the unique morphism into that subalgebra is an inverse to its
inclusion. Thus parameterized induction is also available.

List initiality with parameters gives concatenation
$+\!+:W^2\to W$, membership $\mem:V\times W\to\Two$, and
deletion $\del:V\times W\to W$ by
\begin{align*}
 \nil+\!+w&=w,&
 \cons(y,w)+\!+w'&=\cons(y,w+\!+w'),\\
 \mem(x,\nil)&=\mathbf f,&
 \mem(x,\cons(y,w))&=\epsilon(x,y)\vee_2\mem(x,w),\\
 \del(x,\nil)&=\nil,&
 \del(x,\cons(y,w))&=
 \mathsf{if}\bigl(\epsilon(x,y);\del(x,w),
                         \cons(y,\del(x,w))\bigr).
\end{align*}
In the nil equations the unused variables specify the source of the
constant morphism. These are the algebra-morphism equations of the
corresponding folds. Write $[x]=\cons(x,\nil)$.
Parameterized list induction gives the identities
\begin{align*}
 \mem(x,w+\!+w')&=\mem(x,w)\vee_2\mem(x,w'),\\
 \mem(x,\del(y,w))&=\neg_2\epsilon(x,y)\wedge_2\mem(x,w).
\end{align*}

Give $(W,W)$ the $K$-algebra defined by singleton lists at variables,
ordered concatenation at function and relation symbols and equality,
the identity at negation, concatenation at implication, and deletion
at quantification. Its unique algebra morphism defines
$\FV_{\Tm}:\Tm\to W$ and $\FV_{\Fm}:\Fm\to W$. Explicitly,
\begin{align*}
 \FV_{\Tm}(\eta x)&=[x],&
 \FV_{\Tm}(f_k(s,\vec t))&=\FV_{\Tm}(t_1)+\!+\cdots+\!+\FV_{\Tm}(t_k),\\
 \FV_{\Fm}(\mathsf{eq}(u,t))&=\FV_{\Tm}(u)+\!+\FV_{\Tm}(t),&
 \FV_{\Fm}(r_k(s,\vec t))&=\FV_{\Tm}(t_1)+\!+\cdots+\!+\FV_{\Tm}(t_k),\\
 \FV_{\Fm}(\neg\varphi)&=\FV_{\Fm}(\varphi),&
 \FV_{\Fm}(\varphi\Rightarrow\psi)&=\FV_{\Fm}(\varphi)+\!+\FV_{\Fm}(\psi),\\
 \FV_{\Fm}(\forall x\,\varphi)&=\del(x,\FV_{\Fm}(\varphi)).&&
\end{align*}
An empty concatenation is $\nil$. These are lists of free-variable
occurrences; repetitions are retained.

\begin{definition}\label{def:sentence-object}
The sentence object $s:S\hookrightarrow\Fm$ is the pullback
\[
\begin{tikzcd}
 S\arrow[r,hook,"s"]\arrow[d] & \Fm\arrow[d,"\FV_{\Fm}"]\\
 1\arrow[r,"\nil"'] & W.
\end{tikzcd}
\]
\end{definition}
Negation restricts to $\nu:S\to S$ by the displayed free-variable
equation. Existential quantification is the derived morphism
$\exists=\neg\circ\forall\circ(1_V\times\neg):V\times\Fm\to\Fm$.

\subsection{Admissible substitution}
\label{app:substitution}

We construct substitution of one term for one variable, as required
by the logical schemes below. Put $E=V\times\Tm$, with projections
$x:E\to V$ and $t:E\to\Tm$. For an object $Y$, write
$O(Y)=1+Y$, with injections $\fail_Y:1\to O(Y)$ and
$\ok_Y:Y\to O(Y)$. A morphism $g:Y\to Z$ induces
$O(g)=1_1+g$. For a binary morphism $h:Y^2\to Y$, let
$\widehat h:O(Y)^2\to O(Y)$ return $\ok_Yh$ on the
$Y\times Y$ summand and $\fail_Y$ on the other three summands.
Also let $\widehat\forall:V\times O(\Fm)\to O(\Fm)$ have
restrictions $\fail_{\Fm}$ and $\ok_{\Fm}\forall$.
Distributivity of products over coproducts defines these morphisms.

In $(\mathcal D/E)^2$, use the carriers
\[
 X_T=E\times(\Tm\times\Tm)\longrightarrow E,
 \qquad
 X_F=E\times(\Fm\times O(\Fm))\longrightarrow E.
\]
The first component of each pair records the original expression.
Define a $K_E$-algebra as follows, retaining the same parameter
$(x,t)$ in every operation. At a variable $v$, the pair is
\[
 (\eta v,\mathsf{if}(\epsilon(v,x);t,\eta v)).
\]
At a function symbol, apply its constructor separately to the
original term components and the substituted term components.
At equality and relation symbols, reconstruct the original formula
and apply $\ok_{\Fm}$ to the formula formed from the substituted
term components. At negation and implication, reconstruct the
original formula and apply $O(\neg)$ or
$\widehat{\Rightarrow}$ to the result components.

For quantification, the operation has source
$E\times V\times\Fm\times O(\Fm)$.
Denote its projections by $(x,t,y,\varphi,z)$, with $(x,t)$ the
$E$-coordinate. Define
\[
 C(y,\varphi,x,t)=
 \neg_2\mem(x,\FV_{\Fm}(\varphi))
 \vee_2\neg_2\mem(y,\FV_{\Tm}(t)).
\]
The original component is $\forall y\,\varphi$; the result component is
\begin{equation}\label{eq:quantifier-fold}
 \mathsf{if}\left(\epsilon(y,x);
   \ok_{\Fm}(\forall y\,\varphi),
   \mathsf{if}\bigl(C(y,\varphi,x,t);
          \widehat\forall(y,z),\fail_{\Fm}\bigr)\right).
\end{equation}
In particular, the branch $y=x$ uses the original body and discards
its substitution result.

\begin{proposition}\label{prop:categorical-substitution}
Parameterized initiality determines morphisms
\[
 a_T:\Tm\times E\longrightarrow\Tm,
 \qquad b_F:\Fm\times E\longrightarrow O(\Fm).
\]
Admissibility and formula substitution are defined by the pullback
\[
\begin{tikzcd}
 \Adm\arrow[r,"a_F"]\arrow[d,hook,"j"'] &
 \Fm\arrow[d,hook,"\ok_{\Fm}"]\\
 \Fm\times E\arrow[r,"b_F"'] & O(\Fm).
\end{tikzcd}
\]
The monomorphism $j$ is complemented. These morphisms are uniquely
determined by the preceding algebra and its initial morphism.
\end{proposition}
\begin{proof}
The unique $K_E$-algebra morphism into $(X_T,X_F)$ records an original
and a result component. Composing this algebra morphism with the
projections to the original components gives an endomorphism of the
initial $K_E$-algebra, hence the identity. The result components,
after interchanging the parameter
and syntax factors, are $a_T$ and $b_F$. The pullback defines
$j$ and $a_F$; the coproduct hypotheses make $j$ complemented.
\end{proof}

Put $B=V\times\Fm\times E$ and define
\[
 q=\langle\forall\pi_{V\times\Fm},\pi_E\rangle:B\to\Fm\times E,
 \qquad
 p=\langle\pi_{\Fm},\pi_E\rangle:B\to\Fm\times E.
\]
Let $D_=\hookrightarrow B$ and $D_{\ne}\hookrightarrow B$ be the
two fibers of $\epsilon(y,x)$, and let $D_C=[C=t]$.
Equation~\eqref{eq:quantifier-fold} gives
\begin{equation}\label{eq:admissibility-subobjects}
 q^*\Adm=D_=\;\vee\;
           (D_{\ne}\wedge D_C\wedge p^*\Adm).
\end{equation}
Here meet is pullback intersection and the displayed disjoint join
is the coproduct inclusion into $B$. On $D_=$ the substitution
result is $\forall y\,\varphi$; on the other summand it is
$\forall(y,a_F\widetilde p)$, where
$\widetilde p$ is the induced morphism to $\Adm$.
Thus quantification binds its named variable and prevents capture
of free occurrences of the replacement term.

We write $\varphi[x:=t]$ only on $\Adm$, and $u[x:=t]$ for
$a_T(u,x,t)$. Define
\[
 e_T=\langle\pi_{\Tm},\pi_V,\eta\pi_V\rangle:
          \Tm\times V\to\Tm\times E,
 \qquad
 e_F=\langle\pi_{\Fm},\pi_V,\eta\pi_V\rangle:
          \Fm\times V\to\Fm\times E.
\]
Identity substitution satisfies
\[
 a_Te_T=\pi_{\Tm},\qquad b_Fe_F=\ok_{\Fm}\pi_{\Fm}.
\]
The second equality also states admissibility. For substitution at
an absent variable, use the subobjects
\begin{align*}
 N_T&=[\mem(x,\FV_{\Tm}(u))=f]
                            \hookrightarrow\Tm\times E,\\
 N_F&=[\mem(x,\FV_{\Fm}(\varphi))=f]
                            \hookrightarrow\Fm\times E.
\end{align*}
On $N_T$, $a_T$ is the term projection. On $N_F$, $b_F$ is
$\ok_{\Fm}$ composed with the formula projection, so substitution
is admissible and fixes the formula.

These term and formula statements follow simultaneously by
parameterized induction. Atomic steps use the term equations;
connective steps use the lifted operations. At a quantifier with
equal variable names, \eqref{eq:quantifier-fold} returns the original
formula. For identity substitution with distinct names,
$\FV_{\Tm}(\eta x)=[x]$ makes the capture test true.
For an absent variable with distinct names, the deletion identity
transfers the absence condition to the body. In both cases the
induction hypothesis gives the required successful second branch.

\subsection{Inference rules and deductively closed subobjects}
\label{app:closed-subobjects}

The rule presentation below is the finitary inductive definition of a
formal system~\cite[Section 1.1]{Aczel}, expressed in a slice category.
Let $\mathcal D$ be regular with finite coproducts, and fix an object
$F$. For an integer $N\geq0$, let
\[
 p_n:R_n\longrightarrow F^n,\qquad
 q_n:R_n\longrightarrow F\quad(0\leq n\leq N)
\]
specify rule instances, their ordered premises, and their conclusions.
The nullary rule instances specify axioms.
For an object $q:Q\to F$ of the slice category $\mathcal D/F$,
form the pullback
\[
\begin{tikzcd}
 B_n(Q,q)\arrow[r,"{b_{n,q}}"]\arrow[d,"{a_{n,q}}"'] &
 Q^n\arrow[d,"q^n"]\\
 R_n\arrow[r,"p_n"'] & F^n.
\end{tikzcd}
\]
Put
\[
 H(Q,q)=\left(\coprod_{n=0}^N B_n(Q,q),\kappa_q\right),\qquad
 \kappa_q=[q_na_{n,q}]_{n=0}^N.
\]
For a slice morphism $h:(Q,q)\to(Q',q')$, define $B_n(h)$ by
\[
 a_{n,q'}B_n(h)=a_{n,q},\qquad
 b_{n,q'}B_n(h)=h^n b_{n,q}.
\]
Pullback uniqueness proves the identity and composition laws, so
$H(h)=\coprod_nB_n(h)$ defines an endofunctor of $\mathcal D/F$.

\begin{definition}\label{def:deductively-closed}
A subobject $m:T\hookrightarrow F$ is closed under these rules if
there is a morphism $\alpha:\coprod_nB_n(T,m)\to T$ with
$m\alpha=\kappa_m$. Equivalently, $(T,m)$ admits an $H$-algebra
structure in $\mathcal D/F$. When $F$ and the rule morphisms specify
a logical calculus, such a subobject is a deductively closed theory
of formulas for that calculus.
\end{definition}

The algebra structure is unique because $m$ is monic. Its nullary
part includes all axioms, and its other parts express closure under
the inference rules. Inclusion of closed subobjects is automatically
an algebra morphism, by composing with the target monomorphism.

\begin{theorem}\label{thm:least-theorem-subobject}
With the rule data above, suppose
$\delta:H(P,c)\to(P,c)$ is an initial algebra in $\mathcal D/F$.
The image of its conclusion morphism $c:P\to F$ is the least
subobject of $F$ closed under the rules.
\end{theorem}
\begin{proof}
Factor $c=me$ with $e:P\twoheadrightarrow T$ regular epic and
$m:T\hookrightarrow F$ monic. The endofunctor $H$ preserves
regular epimorphisms: finite products preserve them in a regular
category, each $B_n(e)$ is a pullback of $e^n$, and finite coproducts
preserve coequalizers. Thus $H(e)$ is regular epic.

Writing $|H(P,c)|$ for the underlying object of the slice object,
the equations $c\delta=\kappa_c$ and
$\kappa_mH(e)=\kappa_c$ give the commuting square
\[
\begin{tikzcd}[column sep=large]
 \lvert H(P,c)\rvert\arrow[r,"e\delta"]\arrow[d,two heads,"H(e)"'] &
 T\arrow[d,hook,"m"]\\
 \lvert H(T,m)\rvert\arrow[r,"\kappa_m"']\arrow[ur,dashed,"\alpha"] & F.
\end{tikzcd}
\]
Since $m$ is monic, $e\delta$ coequalizes the kernel pair of $H(e)$.
The regular epimorphism $H(e)$ is its kernel-pair coequalizer, so
there is a unique diagonal $\alpha$. In particular, $m\alpha=\kappa_m$,
so the image is closed. If $u:U\hookrightarrow F$ is another
closed subobject, initiality gives an algebra morphism
$v:(P,c)\to(U,u)$, hence $uv=c$. The image factorization gives
a unique $w:T\to U$ with $uw=m$. This is the required inclusion.
\end{proof}

\subsection{First-order theories}
\label{app:first-order-derivations}

Continue with the syntax structure of
Definition~\ref{def:categorical-syntax} and its constructed substitution.
Put $F=\Fm$. The propositional axiom schemes are those of
H\'ajek and Pudl\'ak~\cite[Section 0.10]{HP}:
\begin{align}
 &\varphi\Rightarrow(\psi\Rightarrow\varphi),
                                                        \label{eq:hilbert-one}\\
 &(\varphi\Rightarrow(\psi\Rightarrow\chi))
   \Rightarrow((\varphi\Rightarrow\psi)
   \Rightarrow(\varphi\Rightarrow\chi)),
                                                        \label{eq:hilbert-two}\\
 &(\neg\psi\Rightarrow\neg\varphi)
   \Rightarrow(\varphi\Rightarrow\psi).
                                                        \label{eq:hilbert-three}
\end{align}
The named-variable quantifier axioms follow
Hirst and Hirst~\cite[Sections 2.11--2.12]{HH}; equality is expressed
by reflexivity and substitution of equal terms in formulas
\cite[Section 3.1]{HH}:
\begin{align}
 &(\forall x\,\varphi)\Rightarrow\varphi[x:=t],
                                                        \label{eq:hilbert-instantiation}\\
 &\forall x(\varphi\Rightarrow\psi)
       \Rightarrow(\varphi\Rightarrow\forall x\,\psi),
                                                        \label{eq:hilbert-distribution}\\
 &t=t,                                                   \label{eq:hilbert-reflexivity}\\
 &(u=t)\Rightarrow
       (\varphi[x:=u]\Rightarrow\varphi[x:=t]).           \label{eq:hilbert-equality}
\end{align}
Each scheme denotes a morphism from its parameter object to $F$.
For the first three these objects are $F^2,F^3,F^2$.
For \eqref{eq:hilbert-instantiation}, use the pullback of
$j:\Adm\hookrightarrow F\times V\times\Tm$ along the
permutation $V\times F\times\Tm\to F\times V\times\Tm$.
For \eqref{eq:hilbert-distribution}, use
$[\mem(x,\FV_{\Fm}(\varphi))=f]
\hookrightarrow V\times F^2$.
For \eqref{eq:hilbert-reflexivity}, use $\Tm$.
For \eqref{eq:hilbert-equality}, intersect the pullbacks of $j$
along
\[
 (x,\varphi,u,t)\longmapsto(\varphi,x,u),\qquad
 (x,\varphi,u,t)\longmapsto(\varphi,x,t)
\]
from $V\times F\times\Tm^2$ to $F\times E$.
Let $L$ be the coproduct of these seven parameter objects and
$\lambda:L\to F$ their copairing.

Fix an axiom subobject $i_A:A\hookrightarrow S$. Define rule objects
and their premise and conclusion morphisms by
\begin{gather*}
 R_0=L+A,\quad p_0:R_0\to1,\quad q_0=[\lambda,si_A],\\
 R_1=V\times F,\quad p_1=\pi_F:R_1\to F,
             \quad q_1=\forall:R_1\to F,\\
 R_2=F^2,\quad p_2=\langle\pi_1,\Rightarrow\rangle:R_2\to F^2,
             \quad q_2=\pi_2:R_2\to F.
\end{gather*}
These are axiom introduction, generalization, and modus ponens.
All nonlogical axioms are sentences. Let $H_A:\mathcal D/F\to
\mathcal D/F$ be the resulting inference endofunctor.

\begin{theorem}\label{thm:first-order-theory}
Suppose $H_A$ has an initial algebra with conclusion morphism
$c_A:P_A\to F$. Its image $T_A\hookrightarrow F$ is the least
subobject containing the logical and nonlogical axioms and closed
under generalization and modus ponens. The image of the pullback
\[
\begin{tikzcd}
 P_{A,S}\arrow[r]\arrow[d,"c_{A,S}"'] & P_A\arrow[d,"c_A"]\\
 S\arrow[r,hook,"s"'] & F
\end{tikzcd}
\]
is $T_{A,S}=S\times_F T_A\hookrightarrow S$.
\end{theorem}
\begin{proof}
Theorem~\ref{thm:least-theorem-subobject} gives the first assertion.
Regular images are stable under pullback, giving the second.
\end{proof}

Suppose these initial derivation algebras exist for every axiom
subobject of $S$. A \emph{sentence theory} is a subobject
$u:U\hookrightarrow S$ through which $c_{U,S}$ factors.
Then $T_{A,S}$ is the least sentence theory containing $A$.
Indeed, axiom introduction gives $A\subseteq T_{A,S}$; since $T_A$
is closed and contains $T_{A,S}$, leastness gives
$T_{T_{A,S}}\subseteq T_A$. Pullback proves closure of $T_{A,S}$.
If $U$ is a sentence theory containing $A$, leastness gives
$T_A\subseteq T_U$, hence
$T_{A,S}\subseteq T_{U,S}\subseteq U$.
It also follows that $A\subseteq A'$ implies
$T_{A,S}\subseteq T_{A',S}$.

In the partial-map category $\Par(\mathcal D,\Mono)$, whose morphisms
are spans with monic left leg composed by pullback
\cite[Section 3]{CockettLack}, the sentence theory is represented by
\[
 t_A=[S\xleftarrow{m_A}T_{A,S}\xrightarrow{m_A}S]:S\rightharpoonup S,
 \qquad m_A:T_{A,S}\hookrightarrow S.
\]
Its restriction is itself. It is the range partial identity of
$c_{A,S}$, meaning the partial identity on that morphism's image.
For a restriction category $\mathcal R$, its restriction splitting
$\Split(\mathcal R)$ has objects $(X,e)$, where $e=\bar e$, and
morphisms $f:(X,e)\rightharpoonup(Y,d)$ satisfying $f=dfe$.
Its identity on $(X,e)$ is $e$; such a morphism is total when
$\bar f=e$. Write $\Tot(\Split(\mathcal R))$ for the subcategory
of total morphisms~\cite[Section 2.3.3]{CockettLack}.
If $\mathcal D$ is this total category, require the theorem inclusion
$m_A$ to admit a partial inverse
$r_A:S\rightharpoonup T_{A,S}$ in $\Split(\mathcal R)$, with
\[
 r_A m_A=1_{T_{A,S}},\qquad m_A r_A=\bar r_A.
\]
Then $\bar r_A:S\rightharpoonup S$ represents the constructed
sentence theory as a restriction idempotent.

\subsection{Propositional theories}
\label{app:propositional-theories}

Let $\mathcal D$ be regular with finite coproducts, fix an atom
object $V\in\operatorname{Ob}(\mathcal D)$, and suppose
$G(Z)=V+Z+Z^2$ has an initial algebra
$S$ with constructors $\theta:V\to S$, $\nu:S\to S$, and
${\Rightarrow}:S^2\to S$. For an axiom subobject
$i_A:A\hookrightarrow S$, let $L=S^2+S^3+S^2$ and let
$\lambda:L\to S$ be given by
\eqref{eq:hilbert-one}--\eqref{eq:hilbert-three}, with $\nu$ for
negation. Use $R_0=L+A$, $q_0=[\lambda,i_A]$, $R_1=0$, and
\[
 R_2=S^2,\qquad
 p_2=\langle\pi_1,\Rightarrow\rangle:S^2\to S^2,
 \qquad q_2=\pi_2:S^2\to S.
\]
If the resulting inference endofunctor has an initial algebra,
Theorem~\ref{thm:least-theorem-subobject} constructs the least
subobject of $S$ containing the axioms and closed under these rules.
Its partial identity belongs to $\Par(\mathcal D,\Mono)$.
In particular, subobjects $i_+:A_+\hookrightarrow V$ and
$i_-:A_-\hookrightarrow V$ specify positive and negative atomic
axioms by the image of
$[\theta i_+,\nu\theta i_-]:A_++A_-\to S$.

\begin{thebibliography}{9}
\bibitem{Aczel}
P.~Aczel, \emph{An introduction to inductive definitions},
in J.~Barwise (ed.), Handbook of Mathematical Logic,
North-Holland, 1977, pp.~739--782.
\href{https://lawrencecpaulson.github.io/papers/Aczel-Inductive-Defs.pdf}{Text}.
\bibitem{CockettLack}
J.~R.~B. Cockett and S.~Lack, \emph{Restriction categories I:
Categories of partial maps}, Theoretical Computer Science
\textbf{270} (2002), pp.~223--259.
\href{https://doi.org/10.1016/S0304-3975(00)00382-0}{doi:10.1016/S0304-3975(00)00382-0}.
\bibitem{HP}
P.~H\'ajek and P.~Pudl\'ak, \emph{Metamathematics of First-Order
Arithmetic}, Perspectives in Logic, Cambridge University Press,
2016 reissue of the 1993 edition.
\href{https://doi.org/10.1017/9781316717271}{doi:10.1017/9781316717271}.
\bibitem{HH}
H.~P.~Hirst and J.~L.~Hirst, \emph{A Primer for Logic and Proof},
2015 edition.
\href{https://www.appstate.edu/~hirstjl/primer/hirst.pdf}{Authors' text}.
\end{thebibliography}
\end{document}
